\chapter{12.4 平移}

\section*{一、选择题（共 8 小题）}

\questionlist
\item
  \begin{questionwithimage}{12.4-1.png}
  如图，$\triangle ABC$ 经过怎样的平移得到 $\triangle DEF$（\quad）
  \optionlist
  \item 把$\triangle ABC$ 向左平移 4 个单位，再向下平移 2 个单位
  \item 把$\triangle ABC$ 向右平移 4 个单位，再向下平移 2 个单位
  \item 把$\triangle ABC$ 向右平移 4 个单位，再向上平移 2 个单位
  \item 把$\triangle ABC$ 向左平移 4 个单位，再向上平移 2 个单位
  \end{enumerate}
  \end{questionwithimage}

\item 在平面内，将一个图形沿某个方向移动一定距离，这样的图形变换为平移，如图，将网格中的三条线段沿网格线的方向（水平或垂直）平移后组成一个首尾依次相接的三角形，至少需要移动（\quad）
  \begin{questionwithimage}{12.4-2.png}
  \optionlist
  \item 12 格
  \item 11 格
  \item 9 格
  \item 8 格
  \end{enumerate}
  \end{questionwithimage}

\item 如图，图案⑥是由①②③④⑤五种基本图形中的两种拼接而成的，这两种基本图形是（\quad）
  \begin{questionwithimage}{12.4-3.png}
  \optionlist
  \item ①⑤
  \item ②③
  \item ④⑥
  \item ②④
  \end{enumerate}
  \end{questionwithimage}

\item 如图所示，半圆 $AB$ 平移到半圆 $CD$ 的位置时所扫过的面积为（\quad）
  \begin{questionwithimage}{12.4-4.png}
  \optionlist
  \item 3
  \item $3+\pi$
  \item 6
  \item $6+\pi$
  \end{enumerate}
  \end{questionwithimage}

\item 如图，在长方形 $ABCD$ 中，$AB=8$，$BC=5$，则图中四个小长方形的周长和为（\quad）
  \begin{questionwithimage}{12.4-5.png}
  \optionlist
  \item 13
  \item 23
  \item 24
  \item 26
  \end{enumerate}
  \end{questionwithimage}

\item 如图，相邻两线段互相垂直，两只蜗牛均同时从点出发到达点，蜗牛甲沿着“$A\rightarrow B\rightarrow C$”路线走，蜗牛乙沿着“$A\rightarrow D\rightarrow E\rightarrow F\rightarrow G\rightarrow H\rightarrow I\rightarrow J\rightarrow C$”的路线走，若他们的爬行速度相同，则先到达点 C 的是（\quad）
  \par\smallskip\noindent
  \begin{minipage}[t]{0.54\linewidth}
  \vspace{0pt}
  \optionlist
  \item 蜗牛甲
  \item 蜗牛乙
  \item 同时到达
  \item 无法确定
  \end{enumerate}
  \end{minipage}
  \hfill
  \begin{minipage}[t]{0.32\linewidth}
  \questionimage{12.4-6.png}
  \end{minipage}
\end{enumerate}

% 以下题目根据截图转录，图形暂未单独拆分为本地图片文件。
\questionlist[7]
\item 如图，将$\triangle ABC$沿$BC$方向平移3 cm得到$\triangle DEF$，若$\triangle ABC$的周长为16 cm，则四边形$ABFD$的周长为（\quad）【图不清楚，可忽略】
  \optionlist
  \item 16 cm
  \item 22 cm
  \item 20 cm
  \item 24 cm
  \end{enumerate}

\item 如图，由$\triangle ABC$平移可得到的三角形有几个（\quad）【图不清楚，可忽略】
  \optionlist
  \item 3个
  \item 5个
  \item 6个
  \item 7个
  \end{enumerate}
\end{enumerate}

\section*{二、填空题（共 6 小题）}

\blanklist[9]

\item 在如图所示的长方体中，可以由线段$AB$平移而得到的线段有\underline{\hspace{3cm}}条。

\centerquestionimage[width=0.55\textwidth]{0.35\textwidth}{12.4-9.png}

\item 如图，在长为$a$米、宽为$b$米的长方形耕地上修筑宽度均为1米的道路（如阴影部分），其余空白部分用于种植草坪，则草坪面积为\underline{\hspace{4cm}}米$^2$。

\centerquestionimage[width=0.55\textwidth]{0.35\textwidth}{12.4-10.png}

\item 在等腰$\triangle ABC$中，$AB=AC$，则$BC$边上的中线、高线和$\angle BAC$的平分线重合于$AD$（如图一）。若将等腰$\triangle ABC$的顶点$A$向右平行移动后，得到$\triangle A'BC$（如图二），那么，此时$BC$边上的中线、$BC$边上的高线和$\angle BA'C$的平分线应依次分别是\underline{\hspace{1.4cm}}、\underline{\hspace{1.4cm}}、\underline{\hspace{1.4cm}}。（填$A'D$、$A'E$、$A'F$）

\centerquestionimage[width=0.6\textwidth]{0.5\textwidth}{12.4-11.png}

\item 如图，把$\angle AOB$沿着直线$MN$平移一定的距离，得到$\angle CPD$。若$\angle AOM=40^\circ$，$\angle DPN=40^\circ$，则$\angle AOB=$\underline{\hspace{2cm}}。

\centerquestionimage[width=0.35\textwidth]{0.55\textwidth}{12.4-12.png}

\item 边长为2的等边$\triangle ABC$与等边$\triangle DEF$互相重合，将$\triangle ABC$沿直线$l$向右平移$m$个单位长度，将$\triangle DEF$向右也平移$n$个单位长度，如图，当$C$、$E$是线段$BF$的三等分点时，$m$的值为\underline{\hspace{2cm}}。

\centerquestionimage[width=0.85\textwidth]{0.35\textwidth}{12.4-13.png}


\item 如图，长方形 $ABCD$ 中，$AB=6$，第一次平移长方形 $ABCD$ 沿 $AB$ 的方向向右平移 5 个单位，得到长方形 $A_1B_1C_1D_1$，第二次平移将长方形 $A_1B_1C_1D_1$ 沿 $A_1B_1$ 的方向向右平移 5 个单位，得到长方形 $A_2B_2C_2D_2$，...，第 $n$ 次平移将长方形 $A_{n-1}B_{n-1}C_{n-1}D_{n-1}$ 沿 $A_{n-1}B_{n-1}$ 的方向平移 5 个单位，得到长方形 $A_nB_nC_nD_n$ ($n \geq 2$)，若 $A_nB_n$ 的长度为 56，则 $n=\underline{\quad }$。

% 插入图片：此处应有第14题图


\vspace{0.2cm}
\centerquestionimage[width=0.90\textwidth]{0.35\textwidth}{12.4-14.png}


\end{enumerate}

\section*{三、解答题（共 5 小题）}


\solutionlist[15]


\item 如图，将 $Rt\triangle ABC$ 沿 $AB$ 方向平移得到 $Rt\triangle DEF$，已知 $BE=6$，$EF=8$，$CG=3$，求阴影部分的面积。

% 第15题图
\rightquestionimage[width=0.55\textwidth]{0.55\textwidth}{12.4-15.png}

\vspace{1cm}

\item 如图，$\triangle A_1B_1C_1$ 是 $\triangle ABC$ 向右平移 4 个单位长度后得到的，且三个顶点的坐标分别为 $A_1(1, 1)$，$B_1(4, 2)$，$C_1(3, 4)$。

(1) 请画出 $\triangle ABC$，并写出点 $A$，$B$，$C$ 的坐标；

(2) 求出 $\triangle AOA_1$ 的面积。

% 第16题图
\rightquestionimage[width=0.75\textwidth]{0.35\textwidth}{12.4-16.png}

\vspace{2cm}

\item 如图，在直角坐标系 $xOy$ 中，$A(-1, 0)$，$B(3, 0)$，将 $A, B$ 同时分别向上平移 2 个单位，再向右平移 1 个单位，得到的对应点分别为 $D, C$，连接 $AD, BC$。

(1) 直接写出点 $C, D$ 的坐标：$C$ \_\_\_\_\_\_，$D$ \_\_\_\_\_\_；

(2) 四边形 $ABCD$ 的面积为 \_\_\_\_\_\_；

(3) $P$ 在线段 $BC$ 上，$\triangle OPD$ 面积是 $\frac{11}{3}$，求 $P$ 点的坐标。

% 第17题图
\rightquestionimage[width=0.75\textwidth]{0.35\textwidth}{12.4-17.png}

\vspace{3cm}


\item 如图，长方形 $ABCD$ 在坐标平面内，点 $A$ 的坐标是 $A(2, 1)$，且边 $AB, CD$ 与 $x$ 轴平行，边 $AD, BC$ 与 $y$ 轴平行，点 $B, C$ 的坐标分别为 $B(a, 1)$，$C(a, c)$，且 $a, c$ 满足关系式：$c = \sqrt{a-6} + \sqrt{6-a} + 3$。


(1) 求 $B, C, D$ 三点的坐标；

(2) 怎样平移，才能使 $A$ 点与原点重合？平移后点 $B, C, D$ 的对应分别为 $B_1C_1D_1$，求四边形 $OB_1C_1D_1$ 的面积；

(3) 平移后在 $x$ 轴上是否存在点 $P$，连接 $PD$，使 $S_{\triangle COP} = S_{\text{四边形} OBCD}$？若存在这样的点 $P$，求出点 $P$ 的坐标；若不存在，试说明理由。

% 第18题图
\rightquestionimage[width=0.85\textwidth]{0.5\textwidth}{12.4-18.png}

\vspace{3cm}

\item 我们约定，若一个三角形（记为 $\triangle A_1$）是由另一个三角形（记为 $\triangle A$）通过一次平移，或绕其任一边的中点旋转 $180^\circ$ 得到的，则称 $\triangle A_1$ 是由 $\triangle A$ 复制的。以下的操作中每一个三角形可以复制一次，复制过程可以一直进行下去。如图 1 是由 $\triangle A$ 复制出 $\triangle A_1$，又由 $\triangle A_1$ 复制出 $\triangle A_2$，再由 $\triangle A_2$ 复制出 $\triangle A_3$，形成了一个大三角形，记作 $\triangle B$。以下各题中的复制均是由 $\triangle A$ 开始的，由复制形成的多边形中的任意两个小三角形（指与 $\triangle A$ 全等的三角形）之间既无缝隙也无重叠。


(1) 图 1 中标出的是一种可能的复制结果，它用到\_\_\_\_\_\_次平移，\_\_\_\_\_\_次旋转。若由复制形成的 $\triangle C$ 的一条边上有 11 个小三角形（指有一条边在该边上的小三角形），则 $\triangle C$ 中含有\_\_\_\_\_\_个小三角形；

(2) 若 $\triangle A$ 是正三角形，你认为通过复制能形成的正多边形是\_\_\_\_\_\_。

(3) 在复制形成四边形的过程中，小明用到了两次平移一次旋转，你能用两次旋转一次平移复制形成一个四边形吗？如果能，请在图 2 的方框内画出草图，并仿照图 1 作出标记；如果不能，请说明理由；

(4) 图 3 是正五边形 $EFGHI$，其中心是 $O$。连接 $O$ 点与各顶点。将其中的一个三角形记为 $\triangle A$，小明认为正五边形 $EFGHI$ 是由复制形成的一种结果，你认为他的说法对吗？请判断并说明理由。

% 第19题图
\rightquestionimage[width=0.7\textwidth]{0.7\textwidth}{12.4-19.png}


\end{enumerate}
